


\section{Modeling}
\label{sec:modeling}



In this section,
we discuss three major modeling components for an \citeeval{} system---retrieval, synthesis, and post-editing.





\begin{table}[t]
    \centering
    \small
    \begin{tabular}{>{\raggedright\arraybackslash\tt}p{0.94\linewidth}<{}}
        \toprule
            Instruction: Write an accurate, engaging, and concise answer for ... \\
            \\
            \vspace{-1em}
            \noindent {\color{brown}<Retrieve for the question>}\\
            Document {[1]}(Title: American Decolonization) ... \\
            Document {[2]}(Title: Decolonization) ... \\
            Document {[3]}(Title: American Revolution) ... \\
            ... \\
            \\
            \vspace{-1em}
            Question: When did  US break away from England? \\
            Answer: {\color{blue} The United States took the first step towards gaining independence ... \tf{\underline{[1][2]}}. The Treaty of Paris was later signed ... \tf{\underline{[3]}}}.\\

        \bottomrule
    \end{tabular}
    \caption{
        An example of our \vani{} method. Different colors represent \ttt{prompt}, \ttt{\color{blue} model generation}, and \ttt{\color{brown} <actions>}. 
        We also provide two in-context demonstrations before the test example. 
    }
    \vspace{-10pt}
    \label{tab:prompt_short}
\end{table}


\subsection{Retrieval}
\label{sec:retrieval}

We explore simple, off-the-shelf retrievers.
We use dense retrievers  for Wikipedia, including GTR~\cite{ni-etal-2022-large} and DPR~\cite{karpukhin-etal-2020-dense}; 
we use 
BM25 for Sphere. %
For each question,
we 
retrieve the top-100 passages. %



\subsection{Synthesis}
\label{sec:synthesis}





We  focus on %
how to
prompt an LLM to 
interact with the retriever, and 
synthesize and cite the evidence (without fine-tuning internal parameters).
One noteworthy challenge is that existing %
LLMs all have limited context window and thus
can only fit a handful of passages.





\paragraph{\textsc{\pvani{}}.}
We simply provide the model with the top-$k$\footnote{We can fit at most $k=3$ for models with 2K window and at most $k=5$ for models with 4K context window.} passages and instruct the model to cite accordingly (Table~\ref{tab:prompt_short}).
We also use in-context learning~\cite{brown2020language} and prepend two demonstrations.
The complete instruction is in Table~\ref{tab:inst_vanilla}.


\paragraph{\textsc{\psumm}/\textsc{\psnippet{}}.}
With a 4K context window, we can at most safely fit $k=5$ passages. As shown in Figure~\ref{fig:ret}, top-5 retrieved passages can only cover 56.8\% percent of the answers in ASQA. 

To tackle this limitation, we propose to provide \textit{summaries} or \textit{snippets} of passages instead of the full text (summaries are abstractive but snippets are spans from passages). 
We acquire summaries and snippets by prompting ChatGPT with instructions (prompts in Table~\ref{tab:prompt_summary} and \ref{tab:prompt_snippet}).\footnote{We also query ChatGPT whether the passage is relevant to the question, and
filter out passages that are ``irrelevant''.}
Then we replace all passages with  summaries/snippets.
Summaries or snippets significantly reduce the passage length, allowing for more passages to fit in:
for ASQA, they reduce passage length by 6$\times$ on average.


Though \summ{}/\snippet{} allows for more retrieved passages, 
they are lossy compressions. 
To alleviate this problem,
we propose \interact{}, an interactive prompting scheme 
to allow the model to check the full text of certain passages. %
At each step, the model can execute one of  three actions: 
(1) ``\ttt{Check: Document [1][2]}'' to check the full text of the corresponding documents;
(2) ``\ttt{Output:}'' to output a statement of the answer;
(3) ``\ttt{End.}'' to end the generation.
\S\ref{app_sec:imp_details} provides more details.






\begin{table}[t]
    \centering
    \small
    \begin{tabular}{>{\raggedright\arraybackslash\tt}p{0.94\linewidth}<{}}
        \toprule
            Instruction: ... \\
            \\
            \vspace{-1em}
            \noindent {\color{brown}<Retrieve for question ``...''>}\\
            Question: When did  US break away from England? \\
        \noindent {\color{blue} Search: Declaration of Independence} \\
        \noindent {\color{brown}<Search the query among the top-100 passages>}\\
        Document [1](Title: ...) ... \\
        \noindent {\color{blue} Output: The United States ... \tf{\underline{[1]}}.} \\
        \noindent {\color{brown}<Remove Document [1] from context>}\\
        \noindent {\color{blue} Search: Treaty of Paris} \\
        \noindent {\color{brown}<Search the query among the top-100 passages>}\\
        Document [3](Title: ...) ... \\
        \noindent {\color{blue} Output: The Treaty of Paris ... \tf{\underline{[3]}}.} \\
        \noindent {\color{brown}<Remove Document [3] from context>}\\
        \noindent {\color{blue}End.}\\
        \bottomrule
    \end{tabular}
    \caption{
        An example of \search{}. 
    }
    \vspace{-15pt}
    \label{tab:prompt_search}
\end{table}


\paragraph{\textsc{\psearch{}}.}
The above methods all display retrieval results at the beginning.
In \search{},
we allow LLMs to call ``search'' during the generation process~\citep{yao2023react,press2022measuring,jiang2023active}.
At each step, the model can execute one of  three actions: 
``\ttt{Search: \{query\}}'' to search among the top-100 passages\footnote{We do not search over the entire corpus because  \ttt{\{query\}} may leave out certain context in the question and searching among the already-retrieved passages gives better results.} by using GTR; the ``\ttt{Output}'' and ``\ttt{End}'' actions are the same as \interact{}.
For each ``\ttt{Search}'' action, we display the best retrieved passage in the context.
The passage is removed after one action to save context space.
Table~\ref{tab:prompt_search} shows an example. %



\paragraph{\textsc{\pclose{}}.}
We also add a simple closed-book baseline, where the model is only prompted with the instruction and the question,
without any retrieved passages provided.
Consequently, this variant does not cite any evidences.

\subsection{Post-editing}
\label{sec:refinement}

In this section we discuss two strategies for refining the  output to further improve its quality.

\paragraph{\textsc{\prerank{}}.}
We randomly sample $n_{\text{sample}}=4$ %
responses for each question, and select the best response using the automatic \textit{citation recall} score. 
we expect \textsc{\prerank{}} to improve the citation quality.

\paragraph{\textsc{\pposthoc{}}.}
For each statement, we find the best matching passage among the top-100 retrieved passages using GTR and cite it.
We combine this with \close{} in our experiments.
